\documentclass[11pt]{article}
% Préambule commun des prétirages v2 (format arXiv : source TeX autonome, pdfLaTeX).
% Moteur : pdfLaTeX (arXiv) → fontenc T1 + inputenc ; XeTeX/LuaTeX (tectonic) → fontspec. Sans ce choix, les
% guillemets « » tapés en dur sortent en « ń ż » sous tectonic (ledger, 2026-11-14).
\usepackage{iftex}
\ifPDFTeX
  \usepackage[T1]{fontenc}
  \usepackage[utf8]{inputenc}
\else
  \usepackage{fontspec}
\fi
\usepackage[french]{babel}
\usepackage{lmodern,microtype}
\usepackage[a4paper,margin=2.5cm]{geometry}
\usepackage{amsmath,amssymb}
\usepackage{graphicx,booktabs,siunitx}
\usepackage{authblk}
\usepackage[numbers,sort&compress]{natbib}
\usepackage[hidelinks]{hyperref}
\usepackage{csquotes}
\sisetup{locale=FR,group-minimum-digits=4}
\renewcommand\Affilfont{\small}
% Étiquette de provenance obligatoire sur chaque nombre (règle du projet) :
\newcommand{\mesure}{\textsuperscript{\textsc{m}}}   % mesuré
\newcommand{\derive}{\textsuperscript{\textsc{d}}}   % dérivé
\newcommand{\suppose}{\textsuperscript{\textsc{s}}}  % supposé

% Ajouts propres à cet article (le préambule commun n'est pas modifié).
\usepackage{pgfplots}
\pgfplotsset{compat=1.16}
\usetikzlibrary{babel}   % les deux-points et points-virgules actifs de babel-french cassent TikZ sans cela

\title{Valider un jumeau numérique de l'offre de soins avant d'en voir les résultats :\\
sept protocoles pré-enregistrés, et ce qui n'y résiste pas}
\author[1]{Équipe hackathon-santé}
\affil[1]{Projet « jumeau numérique de l'offre de soins », dépôt \texttt{hackathon-sante/v2}}
\date{Version du \today}

\begin{document}
\maketitle

\begin{abstract}
Un jumeau numérique de l'offre de soins projette l'effectif « libéral » (secteur privé) de cinq
professions jusqu'en 2050, à toutes les échelles jusqu'à la commune. Les données courent de 2012 à
2026 : aucun backtest n'atteint 2050, et un backtest ne valide pas un levier. Nous avons écrit sept
protocoles, gelés par empreinte SHA-256 avant tout calcul qui les concerne : prévision à dix ans au
plus, cohérence à 2050, effet causal d'un levier, échelle communale, mondes simulés à vérité connue,
prévision probabiliste, et un addendum daté. Nous les avons appliqués tels quels. La régulation de
l'installation infirmière en zone sur-dotée a un effet estimé de $\num{-30.3}$\mesure{}~log-points
au département et de $\num{-32.2}$\mesure{} au sein des bassins de vie. Le test de tendances
parallèles fixé d'avance échoue pourtant deux fois ($p = \num{0.043}$\mesure{} puis
\num{0.025}\mesure{}) : le levier n'est pas validé. Chez les kinés, la pente remonte après
l'annulation de la régulation en 2014 (\num{3.1}\mesure{} points par an). Au backtest, la variante
par défaut ne bat la tendance qu'à l'origine 2021 ; aux origines 2016 et 2018, c'est l'égalité. Dans
des mondes simulés, le moteur est exact sans bruit mais biaisé par le bruit, qu'il lit comme des
entrées. Aucune bande d'incertitude n'est calibrée, et la reproduction des projections DREES à 2050
échoue de 14 à 25\mesure{}~\%. Aucun chiffre 2050 n'est donc citable. Ces échecs, obtenus avant
toute retouche, sont le résultat principal.
\end{abstract}
\noindent\textbf{Mots-clés :} démographie des professions de santé ; validation ; pré-enregistrement ;
différence de différences ; zonage conventionnel ; jumeau numérique.

\medskip
\noindent\textbf{Convention.} Chaque nombre porte sa provenance : \mesure{}~mesuré (par nous, ou par
la source citée à côté), \derive{}~dérivé (calculé à partir d'autres nombres), \suppose{}~supposé
(un choix ou une hypothèse). Tous les nombres viennent des rapports de \texttt{v2/docs/lab/reports/}.

\section{Introduction}
Le projet construit un jumeau numérique de l'offre de soins. Il simule des politiques nationales,
exprimées comme des réglages d'hyperparamètres sourcés, et en mesure l'effet à 2050 sur l'effectif
« libéral » des sages-femmes, chirurgiens-dentistes, pharmaciens, masseurs-kinésithérapeutes et
infirmiers. Ici, « libéral » désigne le secteur privé, lucratif ou non, tout employeur public
exclu. Le décideur visé est un directeur général d'agence régionale de santé (ARS), qui veut voir la
commune.

Trois obstacles rendent la validation difficile. D'abord, les séries du répertoire RPPS couvrent
2012--2026 \citep{drees2026demo}, soit au plus dix ans de recul pour un backtest. L'horizon de
référence, 2050, est 24\derive{}~ans après la dernière donnée. Ensuite, un backtest ne valide pas un
levier : il juge le modèle sur un passé où le paramètre du levier n'a pas bougé. Enfin, le premier
backtest du projet (K1, version~1) a choisi la meilleure de trois variantes \emph{après} avoir vu les
résultats, pour un écart de \num{0.1}\mesure{}~point de MAPE. C'est un biais de
sélection\footnote{Rapport interne \texttt{docs/lab/reports/2026-11-13-k1-backtest.md}, section
« What could make it wrong ».}.

\paragraph{Contributions.}
\begin{itemize}
  \item Six protocoles de validation, un par niveau, chacun avec le test qui le falsifierait, gelés
        par empreinte SHA-256 avant tout calcul (section~\ref{sec:methode}).
  \item Une différence de différences sur le zonage infirmier, avec trois placebos et un test du
        mécanisme du module de régulation du jumeau. Le résultat est rapporté tel quel : levier
        non validé (section~\ref{sec:resultats}).
  \item Une expérience de marche et d'arrêt chez les kinés, dont la prédiction pré-enregistrée tient.
  \item Une étude de cas : les deux revues adversariales du jumeau, qui montrent comment des critères
        d'acceptation peuvent passer par construction (section~\ref{sec:menaces}).
  \item Un code de validation indépendant de l'intérieur du jumeau, applicable à tout estimateur
        qui expose \texttt{fit} et \texttt{predict} (section~\ref{sec:repro}).
\end{itemize}

\section{Travaux existants}
\paragraph{Projections nationales.} La DREES projette les dentistes, sages-femmes et pharmaciens
jusqu'en 2050 avec un modèle par composantes (âge $\times$ sexe $\times$ mode d'exercice), sous un
scénario tendanciel à comportements constants \citep{drees2021dd76}. Les infirmiers ont leur propre
projection à 2050 \citep{er1319}. Ces travaux sont nationaux et ne publient pas leur propre erreur
de prévision. Comparée à nos effectifs 2026, la projection de 2021 sous-estimait de 4 à
7\derive{}~\% en six ans, après correction de l'écart de champ\footnote{Exploration interne
\texttt{docs/lab/explorations/drees-projections-and-apl.md}.}. Nous l'utilisons comme cible de
reproduction, pas comme vérité (niveau~2).

\paragraph{Effet du zonage infirmier.} L'IRDES a estimé l'effet des restrictions à l'installation en
zone sur-dense et des incitations en zone sous-dense, par différence de différences sur le zonage de
2011 et des données 2006--2016 \citep{irdes270}. Il rapporte un écart de densité de $-\num{25.57}$\mesure{}
infirmières libérales pour \num{100000}~habitants en zone sur-dense (relevé par la fiche
\texttt{param-ide-regul-effet-densite} de l'axe~1). Notre niveau~3 reprend ce plan d'expérience sur
une autre série (CartoSanté 2004--2025), une autre unité (le département $\times$ classe de zone) et
une autre mesure (l'effectif en logarithme). Il y ajoute un pré-enregistrement et des placebos.

\paragraph{Tendances parallèles.} Rambachan et Roth \citep{rambachan2023} proposent de borner les
violations des tendances parallèles après traitement par celles observées avant. Leur restriction
$\Delta^{RM}(\bar M)$ borne chaque variation annuelle de la violation après traitement par
$\bar M$ fois la plus grande variation annuelle observée avant. Nous l'utilisons seulement à titre
exploratoire, faute de l'avoir fixée d'avance.

\paragraph{Mondes simulés.} L'observatoire d'un membre de l'équipe \citep{evolutionsoin2026} valide
son modèle sur sept mondes simulés à vérité connue. Il rapporte, sans que nous l'ayons reproduit,
un biais de $+8$ à $+19$\mesure{}~\% sur les seniors et de $-\num{12.4}$\mesure{}~\% à 15~ans quand
les entrées croissent de 3~\% par an. Notre niveau~5 reprend l'idée et transforme ces chiffres en
prédictions.

\paragraph{Prévision probabiliste.} Un score est propre si le prévisionniste maximise son score
espéré en annonçant sa vraie croyance. Le CRPS est propre, et il généralise l'erreur absolue
\citep{gneiting2007jasa}. La démarche recommandée consiste à maximiser la netteté des distributions
prédictives sous contrainte de calibration. La calibration probabiliste revient à l'uniformité des
valeurs PIT, condition nécessaire mais pas suffisante \citep{gneiting2007jrssb}. Pour un
algorithme d'inférence, la calibration par simulation (SBC) vérifie que le rang du paramètre vrai
parmi les tirages a posteriori est uniforme \citep{talts2018sbc}. Ces trois références fondent
notre niveau~6. Il a été proposé par l'état de l'art de l'axe~2\footnote{Exploration interne
\texttt{v2/docs/lab/explorations/2026-11-14-etat-de-lart-methodes.md}, \S~3.6.}.

\section{Données}
\begin{table}[h]
\centering\small
\begin{tabular}{@{}p{3.2cm}p{3.6cm}p{2.3cm}p{5.2cm}@{}}
\toprule
Jeu & Producteur, millésime & Grain & Contrôles \\
\midrule
RPPS \citep{drees2026demo} & DREES, publié le 2026-07-02 & dép.\ $\times$ âge $\times$ sexe $\times$ mode, 2012--2026 & kinés seulement à partir de 2017 (rupture de série) ; infirmiers libéraux jusqu'en 2023 \\
CartoSanté \citep{cartosante2026} & FNPS/SNDS via ARS, extrait le 2026-11-13 & commune $\times$ tranche d'âge, IDE 2004--2025, MK 2005--2025, CD 2006--2025 & la somme des communes égale le département à 1 près dans \num{99.2}\mesure{}~\% des département-années ; Marseille absente \\
Zonage \citep{atlasante2014zonage} & ARS / Atlasanté, métadonnées 2014 & commune, classes 1 à 5 et 9 & \num{36700}\mesure{} communes ; codes des anciennes régions, donc antérieur à 2016 \\
DD76 \citep{drees2021dd76} & DREES, mars 2021 & national, 2030 / 2040 / 2050 & cible du niveau 2 \\
\bottomrule
\end{tabular}
\caption{Données. Toutes sont ouvertes.}
\end{table}

\paragraph{Codage du zonage.} Les classes vont de 1 (très sous-dotée) à 5 (sur-dotée), 9 signifiant
« non défini ». Nous avons vérifié le codage sans regarder les effectifs. Montpellier, Toulon, Brest
et Marseille sont en classe~5 pour les infirmiers ; Paris, Saint-Denis et Lille en classe~3.

\paragraph{Appariement.} Les communes CartoSanté (code géographique 2026) sont appariées au zonage
(code d'environ 2012) par leur code. En 2008, \num{57366}\mesure{} infirmiers sur
\num{57369}\mesure{} sont appariés. Le taux est le même pour les dentistes en 2008 et les kinés en
2011 (plus de \num{99.9}\derive{}~\%).

\section{Méthode}\label{sec:methode}
\subsection{Règles communes}
\begin{itemize}
  \item \textbf{Gel.} Chaque protocole est un rapport écrit avant tout calcul. Une copie immuable est
        rangée dans \texttt{v2/validation/preregistered/}, et son empreinte SHA-256 est inscrite au
        journal des décisions (tableau~\ref{tab:sha}). Tout écart ultérieur est une nouvelle entrée
        datée ; une analyse hors protocole est dite \emph{exploratoire}.
  \item \textbf{Interface.} Un estimateur expose \texttt{fit(t, fenêtre, origine)} et
        \texttt{predict(cible)}, qui rend une série indexée par (profession, territoire). Pour le
        niveau~2, il expose aussi \texttt{trajectory(jusqu'à, scénario)}.
  \item \textbf{Registre.} Une variante n'est confirmatoire que si elle est inscrite et datée
        dans le fichier \texttt{registry.py}, avant son premier passage.
\end{itemize}

\begin{table}[h]
\centering\small
\begin{tabular}{@{}llll@{}}
\toprule
Niveau & Gel & SHA-256 (12 car.) & État \\
\midrule
Plan & 2026-11-14 02:54 & \texttt{0108455d34aa} & --- \\
N1 Prévision & 02:54 & \texttt{eb527ecb06a0} & lancé par l'axe 2 \\
N2 Long terme & 02:54 & \texttt{bc5df3083cc5} & lancé (forme affaiblie) \\
N3 Leviers & 02:54 & \texttt{c37e36457f12} & \textbf{lancé} \\
N4 Échelle fine & 02:54 & \texttt{91f897cc5071} & non lancé \\
N5 Mondes simulés & 03:46 & \texttt{1046545b17c0} & lancé \\
N6 Probabiliste & 03:53 & \texttt{d1cc665cb864} & lancé \\
N7 Addendum & 04:37 & \texttt{06ba47b91647} & lancé \\
\bottomrule
\end{tabular}
\caption{Protocoles gelés. N5, N6 et N7 ont été ajoutés après le lancement de N3, mais avant tout
calcul qui les concerne. Ils s'ajoutent aux autres sans les modifier.}\label{tab:sha}
\end{table}

\subsection{Niveau 1 --- prévision à dix ans au plus}
Les origines sont glissantes. La fenêtre d'ajustement s'élargit à partir de 2012 (kinés : 2017) et
compte au moins 5\suppose{}~ans. Les horizons sont $h \in \{1, 3, 5, 8, 10\}$. La \textbf{métrique
principale}, unique, est la MAPE par département à $h = 5$, sur les départements d'au moins
30\suppose{}~professionnels à l'origine. Elle est moyennée à poids égal sur les cellules
profession $\times$ origine. Soit $\Delta$ l'écart de MAPE entre le candidat et la meilleure des deux
références naïves (persistance, tendance linéaire) :
\[
  \Delta = \mathrm{MAPE}(\text{candidat}) - \min\bigl(\mathrm{MAPE}(\text{persistance}),
  \mathrm{MAPE}(\text{tendance})\bigr).
\]
Le candidat \emph{gagne} si la borne haute de l'IC à 95~\% de $\Delta$ est négative. L'IC vient d'un
bootstrap des départements : 2\,000 tirages, graine \texttt{20261114}, mêmes départements pour
toutes les cellules. La variante par défaut est le stock-flux à taux locaux. Les prévisions de
l'origine 2026 vers 2027, 2028 et 2030 sont \emph{déposées}, avec leur empreinte, avant la
publication du RPPS 2027. C'est le seul test qu'aucune analyse antérieure n'a pu voir.

\subsection{Niveau 2 --- long terme (2050)}
\paragraph{Reproduction.} On impose au moteur les hypothèses du tendanciel DD76 : origine 2020, taux
moyens 2016--2020, entrées tenues au niveau 2016--2020. On compare ensuite des ratios
$g(t) = N(t)/N(2020)$, ce qui neutralise l'écart de champ (+0,8 à +2,5\mesure{}~\% en 2020). La
tolérance est de $\pm 2$\suppose{}~\% en 2030 et de $\pm 3$\suppose{}~\% en 2040 et 2050. On
contrôle aussi la forme : un creux pour les pharmaciens, un plateau pour les sages-femmes. La
tolérance est un choix. Elle reste quatre fois plus petite que l'effet publié d'un quota de
$\pm 20$~\%, soit $\pm 12$ à $\pm 13$\mesure{}~\% en 2050 \citep{drees2021dd76}.

\paragraph{Cas limites.}
\begin{itemize}
  \item \textbf{L1, $h \to 0$} : $\hat N(o) = N(o)$ à $10^{-9}$ près en écart relatif.
  \item \textbf{L2, $h \to \infty$, scénario stationnaire} : l'effectif tend vers
        $N^\star = Q \cdot C$, avec $C$ le nombre de personnes-années d'une cohorte d'un entrant. La
        tolérance est de 1\suppose{}~\%, et deux états initiaux doivent donner le même $N(o+200)$.
  \item \textbf{L3, $h < L$} : un choc de quota n'a aucun effet avant $t_0 + L_p$.
  \item \textbf{L4, conservation} : déjà testée par l'axe~0.
\end{itemize}

\subsection{Niveau 3 --- effet d'un levier par expérience naturelle}
\paragraph{Expérience E1 (infirmiers).} L'avenant~1 à la convention infirmière (2009) n'autorise le
conventionnement en zone « sur-dotée » qu'en remplacement d'un départ. Un nouveau zonage s'applique
à partir de 2012 (avenant~3, arrêté du 25 novembre 2011, et arrêté du 21 décembre 2011 sur la
méthode)\footnote{Dates relevées dans des sources secondaires (questions écrites à l'Assemblée
nationale, presse professionnelle) le 2026-11-14 ; textes primaires non relus.}. La cellule est un
département $d$ $\times$ une classe de zone $z$. On note $Y_{d,z,t}$ la somme des communes appariées
et on pose
\[
  D_{d,t} = \log Y_{d,5,t} - \log Y_{d,3,t}.
\]
Un département entre s'il a au moins 10\suppose{}~infirmiers dans chaque cellule en 2008. Avec deux
classes et des effets fixes département $\times$ année, la différence de différences se réduit à des
moyennes entre départements :
\[
  \beta_k = \overline{D_{\cdot,k} - D_{\cdot,2008}}, \qquad
  \beta = \overline{\textstyle\frac18\sum_{t=2012}^{2019} D_{\cdot,t}
           - \frac15\sum_{t=2004}^{2008} D_{\cdot,t}}.
\]
La période de transition, 2009--2011, est exclue : elle relève d'un zonage de première génération
dont nous n'avons pas la carte. L'estimation s'arrête en 2019, avant le rezonage issu de
l'avenant~6 et la crise sanitaire. Les grappes sont les départements, et l'inférence suit une loi de
Student à $G-1$ degrés de liberté. Cinq tests sont fixés dans cet ordre :
\begin{itemize}
  \item \textbf{T1, tendances parallèles} : $T^2$ de Hotelling sur $(\beta_{2004},\dots,\beta_{2007})$,
        à passer si $p \ge \num{0.10}$\suppose{} ;
  \item \textbf{T2, placebo de date} : un faux traitement en 2006 ;
  \item \textbf{T3, placebo de profession} : les dentistes, dont l'installation n'était pas
        régulée\suppose{} ;
  \item \textbf{T4, placebo de retour à la moyenne} : la zone~4 (très dotée, non régulée) contre la
        zone~3 ;
  \item \textbf{T5, effet}.
\end{itemize}
Un placebo échoue s'il a le même signe que $\beta$ et au moins la moitié\suppose{} de son ampleur.

\paragraph{Test du mécanisme.} Le module \texttt{regulation} du jumeau impose, en zone régulée,
entrées $\le$ sorties. Un plafond parfaitement contraignant tient la zone~5 à son niveau de 2008,
tandis que le contre-factuel croît comme la zone~3 du même département. On compare l'effet que
prédit le module à l'effet observé :
\[
  \beta_{\text{mod}} = \overline{\min\bigl(0, -(\log Y_{\cdot,3,t} - \log Y_{\cdot,3,2008})\bigr)},
  \qquad \theta = \beta_{\text{obs}} / \beta_{\text{mod}} .
\]
Si l'IC de $\theta$ contient 1, le plafond strict survit. S'il tombe dans $]0, 1[$, le module devient
un plafond qui fuit, de paramètre $\theta$. Cette lecture n'est permise que si T1 à T4 passent.

\paragraph{Expérience E2 (kinés).} L'avenant~3 des kinés (arrêté du 10 janvier 2012) instaure la même
régulation. Le Conseil d'État l'annule le 17 mars 2014, pour incompétence : une telle limite relève
de la loi \citep{ce2014mk}. La prédiction pré-enregistrée est double. La pente relative de la zone~5
est négative pendant la régulation (2011--2014). Elle remonte ensuite (2014--2017).

\subsection{Niveau 4 --- échelle fine}
On mesure l'erreur à 5~ans par classe de taille de commune ($1$--$2$, \dots, $>200$), sur les
origines 2014, 2016, 2018 et 2020. L'estimateur par défaut multiplie la projection départementale
par la part communale observée à l'origine. On note $N_{\min}$ la plus petite classe à partir de
laquelle le 80\textsuperscript{e} centile de l'erreur absolue est d'au plus 25\suppose{}~\%. Sous
$N_{\min}$, aucune projection communale n'est affichée sans sa bande. Le bruit de Poisson seul
donne $1{,}28/\sqrt{N} \le 0{,}25$, soit $N \ge 26$\derive{}.

\subsection{Niveau 5 --- mondes simulés à vérité connue}
Le générateur simule 30 départements en âges simples, agrégés en tranches comme le RPPS, avec
20 réplications. Il produit sept mondes : équilibre, pyramide non stationnaire, bruit, cohorte
nombreuse de 55--60 ans, entrées $+3$~\% par an, saut des entrées, report de l'âge de départ.
Quand un monde respecte les hypothèses de l'estimateur, l'erreur nationale doit rester sous
1\suppose{}~\% à 15~ans. Quand il les viole, le signe du biais est prédit d'avance, et doit être
observé dans au moins 16 réplications sur 20. Le générateur ne partage aucun code avec les
estimateurs (section~\ref{sec:menaces}).

\subsection{Niveau 6 --- prévision probabiliste}
N1 juge un point. N6 juge une distribution prédictive : celle que donne l'estimateur s'il rend des
tirages, sinon une distribution empirique construite de la même façon pour tous les estimateurs.
On multiplie la prédiction par $\exp(e)$, où $e$ est tiré parmi les erreurs logarithmiques passées
du même estimateur, de la même profession et au même horizon. Seules comptent les origines dont la
cible ne dépasse pas l'origine courante, avec au moins 20\suppose{} erreurs.
\begin{itemize}
  \item \textbf{Score} : $\mathrm{CRPS}(F, y) = \mathbb{E}|X - y| - \tfrac12\,\mathbb{E}|X - X'|$,
        la forme par échantillon de \citet[éq.~21, au signe près]{gneiting2007jasa}, calculé sur
        $\log N$. Diviser par $y$ rendrait le score impropre ; le passage au logarithme le garde
        propre. Le critère est celui de N1 : bootstrap des départements, borne haute de $\Delta < 0$.
  \item \textbf{Calibration} : la couverture des intervalles centraux doit être de
        $[75 ; 85]$\suppose{}~\% à 80~\% et de $[44 ; 56]$\suppose{}~\% à 50~\%. On juge la
        couverture plutôt qu'un $\chi^2$ sur l'histogramme PIT, parce que les départements d'une
        même origine ne sont pas indépendants. Une bande qui échoue est affichée avec sa couverture
        mesurée, jamais sous son niveau nominal.
  \item \textbf{SBC} : 200 mondes de N5 (monde M3), $L = 99$ tirages a posteriori ou de bootstrap
        paramétrique. On teste l'uniformité des rangs en 20 classes par un $\chi^2$, avec
        $p \ge \num{0.01}$\suppose{} et une correction de Bonferroni.
\end{itemize}
Conséquence dérivée de la grille : à $h = 5$, seule l'origine 2021 a des erreurs passées (celles de
l'origine 2016). Les kinés (série depuis 2017) et les infirmiers libéraux (série jusqu'en 2023) en
sont exclus\derive{}. N6 à $h = 5$ ne jugera donc d'abord que trois professions sur une origine.

\subsection{Addendum N7}
Gelé à 04:37, après le résultat de N3, il ajoute quatre volets.
\begin{itemize}
  \item \textbf{A.} La différence de différences de N3, mais au sein des bassins de vie 2022 (cellule
        bassin $\times$ zone, $Y \ge 5$\suppose{} en 2008, grappes départementales, variance CR1). Mêmes
        tests et seuils ; c'est un \emph{second} test du levier, déclaré comme tel.
  \item \textbf{B.} Le verdict de N1 origine par origine (2016, 2018, 2021), pour tous les estimateurs
        inscrits ; « stable » si le signe de $\Delta$ est le même aux trois origines.
  \item \textbf{C.} Une bande « pente d'erreur » à armes égales :
        $\log N \sim \mathcal N(\log \hat N, (s h)^2)$, $s$ étant la moyenne quadratique des erreurs
        logarithmiques passées divisées par leur horizon ; CRPS gaussien.
  \item \textbf{D.} Une cible ajoutée, l'ER~1319 pour les infirmiers libéraux \citep{er1319}, et la
        réconciliation des dentistes à 2050 par valeurs de Shapley sur trois facteurs (périmètre,
        origine et fenêtre, niveau des entrées), avec un résidu toléré d'au plus 25\suppose{}~\%.
\end{itemize}

\section{Résultats}\label{sec:resultats}
\subsection{Niveau 3, E1 : régulation infirmière}
\begin{table}[h]
\centering\small
\begin{tabular}{@{}lccl@{}}
\toprule
Test & Estimation (log-points $\times 100$) & Seuil & Verdict \\
\midrule
T1 tendances parallèles ($G = 59$) & $F(4, 55) = \num{2.65}$\mesure{}, $p = \num{0.043}$\mesure{} & $p \ge \num{0.10}$ & \textbf{échoue} \\
T2 placebo de date 2006 & $+\num{1.4}$\mesure{} $[\num{-0.7} ; \num{3.4}]$ & IC contient 0 & passe \\
T3 placebo dentistes ($G = 43$) & $+\num{1.7}$\mesure{} $[\num{-1.3} ; \num{4.7}]$ & pas même signe & passe \\
T4 placebo zone~4 ($G = 42$) & $\num{-4.2}$\mesure{} $[\num{-9.6} ; \num{1.3}]$ & $< \frac12|\beta|$ & passe \\
\textbf{T5 effet} & $\mathbf{\num{-30.3}}$\mesure{} $[\num{-34.1} ; \num{-26.4}]$ & IC $< 0$ & détecté \\
\midrule
Pondéré par l'effectif 2008 & $\num{-37.9}$\mesure{} $[\num{-48.9} ; \num{-26.9}]$ & --- & non décisif \\
Transition 2009--2011 & $\num{-1.1}$\mesure{} $[\num{-3.9} ; \num{1.8}]$ & --- & rapporté \\
Zone~1 (incitations) contre zone~3 & $+\num{11.3}$\mesure{} $[\num{5.9} ; \num{16.6}]$ & --- & rapporté \\
\bottomrule
\end{tabular}
\caption{E1, infirmiers, zone sur-dotée (5) contre intermédiaire (3). IC à 95~\%. Source :
\texttt{v2/validation/out/did\_results.json}.}\label{tab:e1}
\end{table}

\textbf{Verdict du protocole : design invalide, levier non validé.} T1 échoue. Le paramètre
$\theta$ n'est donc pas interprété au sens du protocole. Pour mémoire, il vaut
$\num{0.80}$\mesure{} $[\num{0.74} ; \num{0.87}]$ ($\beta_{\text{obs}} = \num{-32.4}$\mesure{},
$\beta_{\text{mod}} = \num{-40.5}$\mesure{}).

\begin{figure}[h]
\centering
\begin{tikzpicture}
\begin{axis}[width=0.9\linewidth, height=6.5cm, xlabel={Année}, ylabel={$\beta_k$ (log-points $\times 100$)},
  xmin=2003.5, xmax=2019.5, ymin=-56, ymax=8, grid=major, legend pos=south west,
  xticklabel style={/pgf/number format/1000 sep={}}, legend style={font=\small}]
\addplot+[mark=*, thick] coordinates {(2004,-3.0) (2005,-2.8) (2006,-2.2) (2007,-2.5) (2008,0) (2009,-1.1)
  (2010,-2.9) (2011,-5.8) (2012,-10.2) (2013,-18.6) (2014,-24.2) (2015,-31.9) (2016,-36.6) (2017,-41.9)
  (2018,-45.6) (2019,-50.1)};
\addlegendentry{IDE, zone 5 contre 3}
\addplot+[mark=square*, dashed] coordinates {(2004,-3.6) (2005,-3.4) (2006,1.1) (2007,1.0) (2008,0) (2009,-1.0)
  (2010,0.2) (2011,-0.4) (2012,-1.6) (2013,-2.7) (2014,-3.9) (2015,-4.6) (2016,-5.2) (2017,-7.3) (2018,-7.3)
  (2019,-8.7)};
\addlegendentry{IDE, zone 4 contre 3 (placebo)}
\addplot+[mark=triangle*, dotted] coordinates {(2006,0.0) (2007,-0.1) (2008,0) (2009,2.1) (2010,2.3) (2011,2.4)
  (2012,1.8) (2013,0.7) (2014,1.7) (2015,1.0) (2016,1.6) (2017,1.9) (2018,1.7) (2019,3.1)};
\addlegendentry{Dentistes, zone IDE 5 contre 3 (placebo)}
\draw[gray] (axis cs:2008.5,-56) -- (axis cs:2008.5,8);
\draw[gray] (axis cs:2011.5,-56) -- (axis cs:2011.5,8);
\end{axis}
\end{tikzpicture}
\caption{Étude d'événement E1, référence 2008\mesure{}. Les traits verticaux encadrent la transition
2009--2011. Les quatre années d'avant sont au même niveau entre elles ($-2{,}2$ à $-3{,}0$) : c'est
2008 qui s'en écarte, ce que T1 rejette. Source : \texttt{did\_e1\_event\_study*.csv}.}
\label{fig:es}
\end{figure}

\paragraph{Lecture.} L'écart entre les années 2004--2007 est faible : de $-\num{2.2}$ à
$-\num{3.0}$\mesure{} points par rapport à 2008. Il prend la forme d'un saut en 2008, pas d'une
tendance. Après 2011, l'écart se creuse de 5 à 8\mesure{}~points par an, sans plateau jusqu'en 2019
(figure~\ref{fig:es}). Ce profil est compatible avec une régulation qui n'agit qu'au flux des
entrées, mais le protocole ne permet pas de conclure.

\paragraph{Analyses exploratoires (écrites après avoir vu T1).} Ces analyses ne valident rien.
\begin{itemize}
  \item En prolongeant la tendance linéaire des années 2004--2008 (\num{+0.64}\mesure{}~point par
        an), l'effet devient $\num{-36.4}$\derive{}. Avant 2009, la zone sur-dotée croissait plus vite :
        la violation joue contre l'effet.
  \item Avec la restriction $\Delta^{RM}(\bar M = 1)$ \citep{rambachan2023}, chaque variation
        annuelle de la violation après 2008 est bornée par la plus grande variation annuelle d'avant,
        \num{2.5}\mesure{}~points. La violation en $t$ vaut donc au plus \num{2.5}\,$(t-2008)$,
        soit un biais d'au plus \num{18.6}\derive{}~points sur la moyenne 2012--2019. En élargissant l'IC de ce biais, on obtient
        $[\num{-52.7} ; \num{-7.9}]$\derive{}. Le signe tient. C'est une approximation (borne
        ajoutée à l'IC), et non l'inférence uniformément valide des auteurs.
  \item L'IRDES trouve un effet de même signe sur une autre série et une autre période
        \citep{irdes270}.
\end{itemize}

\subsection{Niveau 3, E2 : marche et arrêt chez les kinés}
\begin{table}[h]
\centering\small
\begin{tabular}{@{}lccl@{}}
\toprule
Test & Estimation & Prédiction & Verdict \\
\midrule
Tendances parallèles 2005--2010 ($G = 26$) & $F(6, 20) = \num{1.47}$\mesure{}, $p = \num{0.24}$\mesure{} & $p \ge \num{0.10}$ & passe \\
Pente 2011--2014 (pts/an) & $\num{-2.4}$\mesure{} $[\num{-3.5} ; \num{-1.2}]$ & $< 0$ & confirmée \\
Pente 2014--2017 (pts/an) & $+\num{0.7}$\mesure{} $[\num{-0.7} ; \num{2.1}]$ & --- & --- \\
\textbf{Écart des pentes} ($G = 29$) & $\mathbf{+\num{3.1}}$\mesure{} $[\num{1.1} ; \num{5.1}]$ & $> 0$ & \textbf{confirmée} \\
\bottomrule
\end{tabular}
\caption{E2, kinés, zone sur-dotée contre intermédiaire, référence 2011.}\label{tab:e2}
\end{table}
La zone sur-dotée décroche de 7\mesure{}~points pendant les deux ans et demi de régulation (2014 :
$\num{-7.0}$\mesure{}). Après l'annulation, elle cesse de décrocher (2017 : $\num{-4.4}$\mesure{}).
C'est la première prédiction pré-enregistrée sur un levier qui tient. Ce n'est pas encore une
validation du module, car $\theta$ n'a pas été pré-enregistré pour E2.

\subsection{Niveau 1 : premier passage confirmatoire}
La session de l'axe~2 a lancé le banc N1 sur le périmètre privé, après avoir inscrit ses
variantes\footnote{Rapport interne \texttt{v2/docs/lab/reports/2026-11-14-axe2-methodes.md}.}. La
variante par défaut obtient une MAPE de \num{8.68}\mesure{}~\%, contre \num{9.30}\mesure{}~\% pour la
tendance et \num{9.92}\mesure{}~\% pour la persistance. Elle gagne donc :
$\Delta = \num{-0.63}$\mesure{} $[\num{-0.97} ; \num{-0.27}]$. Le seul gain reproduit sur les deux
périmètres est le rétrécissement hiérarchique : $-\num{0.11}$\mesure{} $[\num{-0.21} ;
\num{-0.02}]$. La clause d'instabilité du protocole joue pour les dentistes : la variante par défaut
ne bat la naïve que sur 2\mesure{} origines sur 6. Aucune prévision n'est encore déposée pour 2027.

\subsection{Niveau 3, second test : au sein des bassins de vie (N7-A)}
Sur 135\mesure{} bassins de vie et 60\mesure{} départements, l'effet est de
$\num{-32.2}$\mesure{} $[\num{-37.2} ; \num{-27.3}]$. Le test de tendances échoue de nouveau
($F(4, 56) = \num{3.02}$, $p = \num{0.025}$\mesure{}), ainsi que le placebo de date ($+\num{2.6}$\mesure{}
$[\num{0.1} ; \num{5.1}]$). Les placebos dentistes ($+\num{0.6}$\mesure{}) et zone~4
($\num{-10.3}$\mesure{}, sous la moitié de l'effet) passent. Avant 2008, la zone sur-dotée gagnait
environ 1\mesure{}~point par an sur sa voisine intermédiaire ; après 2011, elle décroche de 6 à
8\mesure{}~points par an. \textbf{Design invalide : le levier IDE reste non validé.} Chez les kinés,
l'écart des pentes est de $+\num{3.3}$\mesure{} $[\num{0.8} ; \num{5.8}]$ points par an, mais le test de
tendances échoue à ce grain ($p = \num{0.059}$\mesure{}).

\subsection{Niveau 1 : stabilité entre origines (N7-B)}
\begin{table}[h]
\centering\small
\begin{tabular}{@{}lccc@{}}
\toprule
Estimateur (périmètre privé) & 2016 & 2018 & 2021 \\
\midrule
Variante par défaut & $+\num{0.10}$ égalité & $\num{-0.22}$ égalité & $\num{-1.22}$ gagne \\
Rétrécissement hiérarchique & $\num{-0.16}$ égalité & $\num{-0.43}$ égalité & $\num{-1.08}$ gagne \\
APC des sorties & $\num{-0.06}$ égalité & $\num{-0.16}$ égalité & $\num{-0.96}$ gagne \\
Rétention 60+ en tendance & $+\num{1.05}$ perd & $+\num{0.17}$ égalité & $\num{-1.25}$ gagne \\
\bottomrule
\end{tabular}
\caption{$\Delta$ MAPE (points) contre la meilleure naïve, $h = 5$, par origine\mesure{}. Extrait.}
\end{table}
Au périmètre privé, la victoire poolée de la variante par défaut repose sur l'origine 2021. Le
classement des estimateurs change presque complètement d'une origine à l'autre (τ de Kendall de
\num{0.07} à \num{0.29}\mesure{}). Sur l'ancien périmètre, la variante par défaut garde le même signe
aux trois origines.

\subsection{Niveau 5 : mondes simulés}
Le moteur est exact dans le monde à l'équilibre (erreur $< 10^{-13}$\mesure{}~\%). Il échoue dans deux
mondes :
\begin{itemize}
  \item \textbf{monde bruité} : 9 réplications sur 20 sous 1~\% d'erreur nationale à 15~ans, et une
        MAPE départementale de \num{10.2}\mesure{}~\% contre un plancher de Poisson de
        \num{5.7}\derive{}~\% ;
  \item \textbf{cohorte nombreuse} : les 60~ans et plus sont \emph{sous}-estimés de
        \num{34.5}\mesure{}~\% à 5~ans, contre la prédiction reprise de l'observatoire.
\end{itemize}
Les mondes à entrées croissantes, à saut d'entrées et à report du départ donnent le signe prédit
(20 réplications sur 20). Avec des entrées en hausse de 3~\% par an, l'ampleur atteint
$\num{-19.8}$\mesure{}~\%, au-delà de la tolérance autour de $-\num{12.4}$~\%. Une même cause explique
les deux échecs : la règle « rétention plafonnée à 1, excédent compté comme entrée » coupe le bruit
vers le bas et transforme en entrées tout excédent, y compris celui d'une cohorte qui passe.

\subsection{Niveau 6 et N7-C : prévision probabiliste}
Avec la bande empirique au même horizon (origine 2021 seulement), la persistance a le meilleur CRPS
(\num{0.0633}\mesure{}) ; la variante par défaut perd ($+\num{0.0098}$ $[\num{0.0037} ; \num{0.0161}]$).
Toutes les bandes sont trop étroites : l'intervalle « 80~\% » couvre de 56 à 70\mesure{}~\% des cas.
Avec la bande « pente d'erreur » (origines 2017--2021), tous les modèles battent la tendance
($\Delta \approx -\num{0.006}$\mesure{}), mais la bande est trop large : l'intervalle « 80~\% »
couvre environ 96\mesure{}~\%. Les tirages propres de l'APC (100 par cellule, moins que les 200
prévus par N6) font jeu égal avec la persistance ($+\num{0.0014}$ $[\num{-0.0033} ; \num{0.0068}]$),
mais leur intervalle « 80~\% » ne couvre que 54\mesure{}~\% des cas. Aucune bande n'est donc publiable sous son niveau nominal. La SBC passe
pour les cinq taux de rétention ($p$ de \num{0.14} à \num{0.94}\mesure{}) et échoue pour les entrées
($p \approx 3 \cdot 10^{-46}$\mesure{}, rang moyen \num{0.22}) : le bruit biaise les entrées vers le
haut, comme dans le monde bruité.

\subsection{Niveau 2 : DREES à 2050 (N2-A, N7-D)}
Faute d'une série ouverte des quotas 2016--2020, la reproduction de DD76 a été lancée avec des
entrées plates (moyenne 2016--2020), version affaiblie du protocole. Elle échoue : en 2050, le ratio
$N(t)/N(2020)$ est inférieur de \num{14.1} (dentistes), \num{15.0} (pharmaciens) et \num{17.4}\mesure{}~\%
(sages-femmes) à celui de DD76 ; les formes (creux, plateau) sont respectées. Une variante à entrées
par la filière dépasse DD76 de 14 à 15\mesure{}~\% (analyse exploratoire, qui utilise des inscrits
connus après 2020). L'ER~1319 part du même effectif 2021 que nous (\num{98631}\mesure{}) ; en 2050 nous
sommes \num{25.2}\mesure{}~\% en dessous, mais l'observé 2023 est déjà \num{5.9}\mesure{}~\% sous
l'ER~1319. Pour les dentistes, notre cible à entrées récentes (\num{81491}\mesure{} en 2050) dépasse
DD76 (\num{61800}) de \num{19691}. Les valeurs de Shapley attribuent \num{13056}\mesure{} à l'origine et
à la fenêtre, \num{11658}\mesure{} aux entrées récentes et $-\num{1877}$\mesure{} au périmètre ; le résidu
(16\derive{}~\%) passe le seuil. La variante par défaut, à entrées moyennes, donne
\num{61964}\mesure{}, proche de DD76.

\subsection{Niveau 4}
Il n'est pas lancé. Son code a été vérifié sur données synthétiques : avec une erreur de Poisson pure,
$N_{\min}$ tombe dans la classe 21--50\mesure{}, en accord avec la dérivation ($N \ge 26$).

\section{Validation et menaces à la validité}\label{sec:menaces}
\subsection{Cas d'étude : des tests vrais par construction}
Le jumeau a subi deux revues adversariales\footnote{Rapports internes
\texttt{2026-11-14-revue-adversariale-axe0.md} et \texttt{-v2.md}.}. Neuf agents en lecture seule
les ont menées selon trois angles. Pour chaque angle, les deux constats les plus graves ont été
soumis à un agent chargé de les réfuter en exécutant le code.
\begin{itemize}
  \item \textbf{Première revue :} 6\mesure{} constats soumis à réfutation, 6 confirmés, dont
        4 bloquants. Le critère « plug-in » écrivait le module factice dans un dossier temporaire :
        son contrôle « aucun fichier modifié » ne pouvait pas échouer. Le critère de conservation
        comparait le registre des flux à une variation qu'il stockait lui-même, si bien que son
        résidu de $3 \cdot 10^{-11}$ ne mesurait que l'arrondi. L'agrégation comparait les communes
        à des regroupements des mêmes lignes ; face au RPPS, l'écart est de
        $-\num{5.4}$\mesure{}~\% pour les kinés en 2025.
  \item \textbf{Seconde revue :} 6\mesure{} constats confirmés, soit 4 défauts bloquants
        distincts. Le test du levier retraite comparait le jumeau à la version~1, qui portait le
        même masque fautif : il ne pouvait pas voir le défaut (+41\mesure{}~\% d'effet à 2050 pour un
        report de 2~ans). Une correction en avait aussi introduit un nouveau : un scénario passé à
        \texttt{trajectory} remplaçait le scénario de base ($-\num{15.2}$\mesure{}~\%).
\end{itemize}
\textbf{Leçon pour nos protocoles.} Un test qui compare un objet à une copie de lui-même, ou qui
mesure une identité, passe par construction. Nos protocoles s'en gardent de quatre façons :
\begin{itemize}
  \item N5 interdit au générateur de partager du code ou une hypothèse avec l'estimateur ;
  \item N2 compare à une source extérieure, DD76, et pas à la version~1 ;
  \item N3 ajoute des placebos qui \emph{peuvent} échouer ;
  \item la règle de mutation de l'axe~0 (réintroduire le défaut et vérifier que le test échoue)
        est la contrepartie, dans le code, du pré-enregistrement.
\end{itemize}

\subsection{Menaces propres à nos résultats}
\begin{itemize}
  \item \textbf{Date du zonage.} Le jeu Atlasanté porte des codes de régions antérieurs à 2016 et
        des métadonnées de 2014 : il est \emph{probablement} le zonage de 2012. S'il s'agissait
        d'une révision de 2014 fondée sur des densités postérieures à 2012, le traitement serait en
        partie choisi sur le résultat. E1 et E2 tomberaient ensemble.
  \item \textbf{Un verdict qui tient à un seuil.} Le test donne $p = \num{0.043}$ pour un seuil de
        \num{0.10}. Ce seuil était fixé d'avance, et nous ne le changeons pas. Un second test, par
        exemple $\Delta^{RM}$ avec $\bar M$ fixé d'avance, devra être pré-enregistré et déclaré
        comme tel.
  \item \textbf{Retour à la moyenne.} La zone~4 décroche aussi, de $-\num{4.2}$\mesure{}~points en
        moyenne. Jusqu'à environ un septième\derive{} de l'effet IDE peut venir de la convergence
        des zones denses.
  \item \textbf{Dates d'effet.} Elles viennent de sources secondaires. L'absence de régulation
        dentaire avant 2025 et la date du rezonage infirmier sont supposées.
  \item \textbf{Contamination de N1.} Les cibles 2023 et 2026 avaient déjà été vues par K1. Seul le
        dépôt de 2027 est propre. Le bootstrap par département ignore aussi la corrélation entre
        origines, si bien que l'IC est probablement trop étroit\suppose{}.
  \item \textbf{Ordre de gel.} N5 a été gelé après le lancement de N3. Les deux niveaux ne partagent
        ni données ni seuil.
  \item \textbf{Seconds tests.} N7-A est un second essai sur le même levier : la protection du
        pré-enregistrement est entamée. Il échoue de toute façon.
  \item \textbf{Forme affaiblie de N2.} Faute de série de quotas, les entrées de DD76 ne sont pas
        reproduites ; l'échec de N2-A mêle l'écart de moteur et l'écart d'hypothèse.
  \item \textbf{Seuils choisis.} Les seuils de 3~\% (N2), 25~\% (N4) et $p \ge 0{,}10$ (N3) sont des
        choix argumentés, pas des dérivations.
\end{itemize}

\section{Reproductibilité}\label{sec:repro}
Toutes les commandes se lancent depuis la racine du dépôt. \texttt{scipy} est ajouté à la volée,
sans toucher au \texttt{pyproject.toml} partagé.
\begin{verbatim}
# N3, N3 exploratoire, auto-test, N7-A, N7-B, N6 et N7-C, N5 et SBC, N2
PYTHONPATH=v2 uv run --with scipy python -m validation.did
PYTHONPATH=v2 uv run --with scipy python -m validation.did_explore
PYTHONPATH=v2 uv run --with scipy python -m validation.selftest
PYTHONPATH=v2 uv run --with scipy python -m validation.did_commune
PYTHONPATH=v2 uv run --with scipy python -m validation.stability prive
PYTHONPATH=v2 uv run --with scipy python -m validation.probabilistic
PYTHONPATH=v2 uv run --with scipy python -m validation.worlds all
PYTHONPATH=v2 uv run --with scipy python -m validation.dd76
# empreintes des protocoles gelés
shasum -a 256 v2/validation/preregistered/*.md
\end{verbatim}
\begin{itemize}
  \item \textbf{Graines} : \texttt{20261114} pour tous les bootstraps ; \texttt{20261114}~+~$r$ pour
        la réplication $r$ des mondes simulés.
  \item \textbf{Données} :
        \begin{itemize}
          \item \texttt{data/processed/stock\_commune.parquet} (CartoSanté, extrait le 2026-11-13) ;
          \item \texttt{v2/validation/raw/zonage\_conv\_c987.geojson} (Atlasanté, extrait le
                2026-11-14 ; libre réutilisation en citant la source et la date) ;
          \item \texttt{data/processed/stock.parquet} (RPPS, publié le 2026-07-02).
        \end{itemize}
  \item \textbf{Sorties} : \texttt{v2/validation/out/}.
  \item \textbf{Code} : les modules Python du dossier \texttt{v2/validation/} (un par niveau,
        plus le registre des estimateurs).
  \item \textbf{Autres données} : tables communales de l'axe~4
        (\texttt{v2/data/processed/}) ; projections DREES des infirmières 2024 (data.drees, export
        CSV, extrait le 2026-11-14).
  \item \textbf{Version} : le dépôt n'était pas encore versionné au moment des calculs ; les
        empreintes des protocoles tiennent lieu de version.
\end{itemize}

\section{Conclusion}
La validation d'un jumeau à 2050 ne peut pas être un seul backtest. Nous l'avons découpée en sept
protocoles, gelés avant leurs résultats et appliqués tels quels. Presque tous les tests qui pouvaient
échouer ont échoué, et c'est l'information principale :
\begin{itemize}
  \item le levier IDE a un effet grand et constant, mais deux designs fixés d'avance refusent de le
        valider ;
  \item la précision par département n'est acquise que sur une origine ;
  \item le moteur par défaut confond bruit, passage de cohorte et entrées ;
  \item aucune bande n'est calibrée ;
  \item à 2050, l'hypothèse d'entrées domine tout, de $-15$ à $+15$\mesure{}~\% autour de DD76.
\end{itemize}

\paragraph{Ce que cela change.}
\begin{itemize}
  \item Aucun chiffre 2050 n'est citable, et aucune bande n'est affichée sous son niveau nominal.
  \item Les axes 0 et 2 doivent séparer les entrées de la rétention.
  \item L'axe~5 doit vérifier la calibration de sa composante (b), linéaire en $h$.
  \item Le niveau d'entrées de référence (moyenne de fenêtre ou années récentes) est une décision qui
        revient à l'équipe : il fait passer les dentistes en 2050 de \num{62000} à \num{81000}\mesure{}.
\end{itemize}

\paragraph{Questions ouvertes, avec leur coût (temps de travail estimé).}
\begin{itemize}
  \item Bande en $\sqrt h$ pré-enregistrée, puis test de couverture : environ 30~min.
  \item Troisième test du levier IDE (contrôle synthétique par bassin de vie) : environ 1~h.
  \item Série ouverte des quotas 2016--2020 pour une reproduction fidèle de DD76 : environ 45~min.
  \item Dater le jeu Atlasanté contre un arrêté ARS de 2012 : environ 30~min.
  \item Cas limites L1--L3 sur le jumeau et niveau~4 : environ 30~min.
  \item Dépôt des prévisions 2027--2030 : environ 5~min.
\end{itemize}

\bibliographystyle{plainnat}
\bibliography{refs}
\end{document}
